\section{Resumen}

Los mercados cripto híbridos combinan mecanismos centralizados de negociación con
creadores automáticos de mercado descentralizados (AMM). El arbitraje puede reducir
discrepancias entre precios, mientras que las restricciones de ejecución, las
decisiones heterogéneas y los cambios endógenos de liquidez pueden modificar la
dinámica de ajuste. Esta tesis propone un marco de dos niveles: (i) un mapa discreto
reducido con variables observables en el simulador y (ii) un micromundo estocástico con
referencia CEX externa, AMM de producto constante y cuatro roles representativos
provisionales. El MVP reproducible contiene comprobaciones mecánicas deterministas y
un modelo de selección emparejado de 300 semillas. En los escenarios configurados, una
mayor capacidad de arbitraje reduce el error máximo de seguimiento, mientras que el
efecto del LP adaptativo cambia según la dirección del shock. Con semillas y magnitudes
de shock reservadas, la etiqueta de estabilidad de Schur del mapa estimado coincide con
el resultado mayoritario de contracción en solo 3 de 8 grupos. Estos resultados
sintéticos de selección no validan el modelo frente a mercados reales ni convierten los
coeficientes del mapa en parámetros empíricos. Aún se requiere revisar el diseño de
agentes y ejecutar experimentos confirmatorios más amplios.

\textbf{Palabras clave:} seguimiento CEX/DEX; AMM; arbitraje; liquidez adaptativa;
simulación multiagente; mapa discreto; reproducibilidad.
